\section{Introduction}

\begin{frame}
  Introduction
\end{frame}

\begin{frame}
  \frametitle{What's new?}
  \begin{itemize}
    \item Queimadas method:
      \begin{itemize}
        % \item Outliers tables~\hyperlink{outliers_and_cis}{\beamergotobutton{}}
        % \item Confidence intervals
        %   ~\hyperlink{outliers_and_cis}{\beamergotobutton{}}
        % \item Use of a month-by-month regression and forecast
        %   \hyperlink{quiemadas_1_month}{\beamergotobutton{}}
        % \item Forecasts' differentes between Queimadas 12 vs. month-by-month
        %   ~\hyperlink{12_vs_1_month_differences}{\beamergotobutton{}}
        % \item Brazilian fire foci aggregated by year and month (facets) with
        %   linear regression of annual peaks
        %   ~\hyperlink{lm_peak_facets}{\beamergotobutton{}}
        % \item BFast analysis
        %   ~\hyperlink{bfast}{\beamergotobutton{}}
        % \item BFast break of anual top values for AQUA M-T time series
        %   ~\hyperlink{bfast_aqua}{\beamergotobutton{}}
        \item Updated \textit{Queimadas} database to include data from 1998.
            ~\hyperlink{lm_peak}{\beamergotobutton{}}
        \item BFast:
        \begin{itemize}
          \item AQUA\_M-T still breaks on 2007-11-15.
            ~\hyperlink{bfast_aqua}{\beamergotobutton{}}
          \item New break on NOAA-12 on 2005-11-15.
            ~\hyperlink{bfast_noaa}{\beamergotobutton{}}
        \end{itemize}
      \end{itemize}
  \end{itemize}
\end{frame}

\begin{frame}
  \frametitle{What is a reference satellite?}
  \begin{columns}
    \begin{column}{0.5\textwidth}
       The reference satellite is~\cite{setzer2020}:
      \begin{itemize}
        \item It's the satellite that daily senses fire foci used to build time
          series. 
        \item NOAA-12 (AVHRR sensor, late afternoon pass) from 1998-06-01 to
          2002-07-03.
        \item Then, AQUA\_M-T (MODIS sensor, early afternoon pass).
        \item The reference satellite sennses only a fraction of the total fire
          foci.
        \item Evenso, this allows trend analysis.
        \item Next reference satellite, NPP-SUOMI (VIIRS sensor).
      \end{itemize}
    \end{column}
    \begin{column}{0.5\textwidth}
      \begin{figure}
        \centering
        \includegraphics[scale=0.5]{logos/logo_queimadas.png}
      \end{figure}
    \end{column}
  \end{columns}
\end{frame}

\begin{frame}
  \frametitle{What would be the next reference satellite?}
  \begin{columns}
    \begin{column}{0.5\textwidth}
      \begin{itemize}
        \item What would be the next reference
          satellite~\cite{bdqueimadas2026,setzer2019banco}?
      \end{itemize}
    \end{column}
    \begin{column}{0.5\textwidth}
      \begin{figure}
        \centering
        \includegraphics[scale=0.5]{logos/logo_queimadas.png}
      \end{figure}
    \end{column}
  \end{columns}
\end{frame}

\begin{frame}
  \frametitle{Online resources}
  \begin{itemize}
    \item Queimadas \href{https://terrabrasilis.dpi.inpe.br/queimadas/portal/}
      {portal}.
    \item Queimadas
      \href{https://terrabrasilis.dpi.inpe.br/queimadas/portal/pages/secao_informacoes/faq/index.html}
      {FAQ}.
    \item Source code
      \begin{itemize}
      \item Code repository available 
        \href{https://github.com/albhasan/queimadas}{here}.
      \item Slides available \href{https://github.com/albhasan/slides/queimadas}
        {here}.
      \end{itemize}
  \end{itemize}
\end{frame}
